% Written By Enoch Yu
% Downloaded from archive.enochyu.com

\documentclass{article}

\usepackage{amsmath}
\usepackage{amsthm}
\usepackage{amssymb}
\usepackage{gensymb}

\usepackage[font=small,labelfont=bf]{caption}
\usepackage{float}
\usepackage[margin=1in]{geometry}
\usepackage{enumitem}

\usepackage{fancyhdr}
\pagestyle{fancy}
\fancyhead[L]{Ramanujan Summation}
\fancyhead[R]{Enoch Yu}

\newtheorem{theorem}{Theorem}[section]
\newtheorem*{theorem*}{Theorem}
\newtheorem{lemma}[theorem]{Lemma}
\newtheorem{sublemma}{Lemma}[section]
\newtheorem{proposition}{Proposition}
\newtheorem{corollary}{Corollary}[theorem]
\newenvironment{solution}{\begin{trivlist}\item[]{\bf Solution}}{\qed \end{trivlist}}

\title{Ramanujan Summation}
\author{Enoch Yu}
\date{1 May 2025}

\begin{document}

\begin{theorem*}
The sum of all natural numbers are $-\frac{1}{12}$, or
\[
\sum_{n=1}^{\infty}n=-\frac{1}{12}.
\]
\end{theorem*}
\begin{proof}
Let $S_1,S_2\text{ and }S_3$ be following sequences.
\begin{align*}
    S_1 &= 1 - 1 + 1 - 1 + 1 - 1 + \cdots \\
    S_2 &= 1 - 2 + 3 - 4 + 5 - 6 + \cdots \\
    S_3 &= 1 + 2 + 3 + 4 + 5 + 6 + \cdots
\end{align*}
Despite the fact that $S_1$ is a divergent sequence,
the following property may apply.
\begin{align*}
    S_1 &= 1 - 1 + 1 - 1 + 1 - 1 + \cdots \\
    S_1 &= \,\,\,\,\,\text{ }\text{ }\text{ }
        1 - 1 + 1 - 1 + 1 - \cdots \\
    2S_1 &= 1 \\
    \therefore
    S_1 &= \frac{1}{2}
\end{align*}
Similarly, $S_2$ may be manipulated.
\begin{align*}
    S_2 &= 1 - 2 + 3 - 4 + 5 - 6 + \cdots \\
    S_2 &= \,\,\,\text{ }\text{ }\text{ }\text{ }
        1 - 2 + 3 - 4 + 5 - \cdots \\
    2S_2 &= 1 - 1 + 1 - 1 + 1 - 1 + \cdots \\
    2S_1 &= S_1 \\
    2S_1 &= \frac{1}{2} \\
    \therefore
    S_2 &= \frac{1}{4}
\end{align*}
Lastly, the relationship between $S_2$
and $S_3$ may be investigated.
\begin{align*}
    S_3 &= 1 + 2 + 3 + 4 + 5 + 6 + \cdots \\
    S_2 &= 1 - 2 + 3 - 4 + 5 - 6 + \cdots \\
    S_3 - S_2 &= 4 + 8 + 1 2 + 16 + \cdots \\
    &= 4(1 + 2 + 3 + 4 + \cdots) \\
    &= 4S_3 \\
    3S_3 &= -S_2 \\
    3S_3 &= -\frac{1}{4} \\
    \therefore
    S_3& = -\frac{1}{12}
\end{align*}
\end{proof}
\end{document}

